\chapter{Conclusion}

The wide range of applicability of \cevns\ in many disciplines in physics is of immense motivation to continue to further our knowledge of the subject. Of particular importance and the main motivation for this study, is the immense role that nuclear structure physics has played in furthering our understanding of the \cevns\ cross-section. The weak form factor, being the last piece of physics for predictive capabilities of \cevns , was calculated based on relativistic mean-field potentials. Using the given model-dependent, complex mean-field potentials for the interactions of the individual nucleons, the point-nuclear density distributions were constructed. It was shown that the structure of the point-nuclear density distributions near the origin was largely controlled by the energy scheme that the model predicted, namely the $l=0$ occupied energy levels. Furthermore, due to the energy differences of the order of $\leq 1$ MeV, the spectra for the same nuclei could differ between models. The structural differences of the density distribution between models were qualitatively described in terms of the surface diffuseness and half-density radius to describe the form factor---precise at large values of momentum transfer. It was shown that the half-density radius controls the form factor oscillations and the surface diffuseness was responsible for the exponential decay. Moreover, the exponential fall-off of the form factor could also be described by the size of the mean-squared radius to first order, where the larger radius of the density distribution would reflect larger exponential decay in the form factor expansion. The radius of the density distributions were also calculated for the various nuclei studied to show the model dependence and validate the predicted neutron distributions by comparing our results with the known proton radii. The neutron-rich nuclei with a larger radius gave much larger corrections to the \cevns\ cross-section. Due to the extremely large corrections up to $\sim 40 \% $ for Lead-208, the precise measurement of the weak form factors is necessary. Using the predicted weak form factors, the differential and total cross-sections were calculated. The relationship between the enhancement to the cross-section due to the number of neutrons and the threshold for the nuclear recoil, a unique experimental limitation, was shown. Finally, the model-independent analysis that was formulated for the extraction of the radius of the density distributions and weak radius, as well as the precision needed to constrain the models could guide experiments that probe the weak interaction through \cevns\ in the future.  